%\documentclass[12pt,oneside]{amsart}
\documentclass{scrartcl}

%\documentclass{article}

\usepackage{amsthm,amsmath,amssymb}
\usepackage[numbers]{natbib}

%\usepackage[colorlinks]{hyperref}
\usepackage{hyperref}
\hypersetup{
    colorlinks,%
    citecolor=black,%
    filecolor=black,%
    linkcolor=black,%
    urlcolor=black
}
\usepackage{hypernat}
\usepackage[top=2.5cm, bottom=2.5cm, left=2.8cm,
right=2.8cm]{geometry}
\usepackage{graphicx} 
\usepackage{pdfpages}

%\setlength{\parskip}{0pt}
%\setlength{\parsep}{0pt}
%\setlength{\headsep}{0pt}
%\setlength{\topskip}{0pt}
%\setlength{\topmargin}{0pt}
%\setlength{\topsep}{0pt}
%\setlength{\partopsep}{0pt}

\linespread{1}

%\usepackage{marvosym}

%\textwidth = 6.5 in \textheight = 9.2 in \oddsidemargin = 0.0 in
%\evensidemargin = 0.0 in \topmargin = -0.0 in \headheight = 0.0 in
%\headsep = 0.0in \parskip = 0.0in \parindent = 0.0in
%\def\baselinestretch{0.871}


\newtheorem{theorem}{Theorem}[section]
\newtheorem{proposition}[theorem]{Proposition}
\newtheorem{problem}[theorem]{Problem}
\newtheorem{fact}[theorem]{Fact}
\newtheorem{lemma}[theorem]{Lemma}
\newtheorem{defn}[theorem]{Definition}
\newtheorem{corollary}[theorem]{Corollary}
\newtheorem{remark}{Remark}[section]
\newtheorem{example}[theorem]{Example}

\newcommand{\E}{{\mathbb E}}
\newcommand{\se}{{\mathcal{E}}}
\newcommand{\R}{{\mathbb R}}
\renewcommand{\P}{{\mathbb P}}
\newcommand{\ac}{{\mathcal{A}}}
\newcommand{\A}{{\mathcal{A}}}
\newcommand{\B}{{\mathcal{B}}}
\newcommand{\C}{{\mathcal{C}}}
\newcommand{\M}{{\mathcal{M}}}
\newcommand{\Mf}{{\mathfrak{M}}}
\newcommand{\T}{{\mathcal{T}}}
\newcommand{\Z}{{\mathcal{Z}}}
\newcommand{\K}{{\mathcal{K}}}
\newcommand{\lc}{{\mathcal{L}}}
\newcommand{\Ps}{{\mathcal{P}}}
\newcommand{\G}{{\mathcal{G}}}
\newcommand{\pa}{{\mathring{p}}}
\newcommand{\F}{{\cal F}}
\newcommand{\samp}{{\mathcal{X}}}
\newcommand{\X}{{\mathcal{X}}}
\newcommand{\hi}{{\hat{\Phi}}}
\newcommand{\data}{{\text{data}}}

\newcounter{rcnt}[section]
\renewcommand{\thercnt}{(\roman{rcnt})}

\def\qt#1{\qquad\text{#1}}

\def\argmin{\mathop{\rm argmin}}
\def\argmax{\mathop{\rm argmax}}


\setlength{\parskip}{1.4 \medskipamount}
%\sloppy
%\linespread{1.3}

\begin{document}

\title{\textbf{Diagnostic Classification Models: Rigorous Mathematical 
Formulation of DINA, DINO, rDINA, and rDINO Models 
with Prior Specifications}}
\author{\textbf{Jonathon A. Pedroza \& Boting Zhang}}
\date{\today}

\maketitle


\begin{abstract}
Diagnostic classification models (DCMs) infer a respondent's latent mastery profile from 
multivariate binary item responses. This project formulates and compares conjunctive and 
disjunctive DCMs for a five-attribute probability-and-statistics assessment. 
Let $Y_{ij}\in\{0,1\}$ denote the response of respondent $i$ to item $j$, and 
let $Z_i\in\{1,\ldots,2^K\}$ index the respondent's latent attribute class induced 
by a prespecified $Q$-matrix. Conditional on $Z_i=l$, responses are modeled as independent 
Bernoulli variables with success probabilities determined by an ideal-response indicator: 
the DINA model uses an AND rule, whereas the DINO model uses an OR rule. The unrestricted 
models estimate item-specific slipping and guessing parameters; the restricted rDINA and 
rDINO models replace these by a common error parameter. Marginalizing over the $2^K$ latent 
classes yields a finite-mixture likelihood with unknown class proportions. We compare 
maximum likelihood estimation, maximum a posteriori estimation, and Bayesian posterior 
inference by Hamiltonian Monte Carlo under uniform and weakly informative Beta priors on 
item-error parameters. The analysis evaluates the effect of structural restrictions and 
regularization on numerical stability, latent-class recovery, posterior uncertainty, 
and predictive adequacy using posterior predictive replications, item-level discrepancies, 
root mean squared error. In particular, the study treats 
prior specification as a statistical regularization device for mitigating boundary 
solutions in sparse latent-class models, rather than solely as a subjective Bayesian input. 
The resulting workflow provides a reproducible framework for estimating, diagnosing, and 
comparing cognitively interpretable latent-variable models for educational assessment data.
\end{abstract}

\end{document}